\documentclass[11pt]{article}

\usepackage[margin=1in]{geometry}
\usepackage{amsmath}
\usepackage{amssymb}
\usepackage{hyperref}

\title{Stage 3H Physics Note: CCZ4 RHS Block Decomposition}
\author{GL-with-AI Project}
\date{\today}

\begin{document}
\maketitle

\section{Scope}

Stage 3H is a planning note for the future \(4+1\) CCZ4 / SO(3)
modified-cartoon right-hand side. It does not implement C++ source terms, add
variables, derive the full nonlinear RHS, or prove physical correctness. Its
purpose is to split the future RHS into separately checkable blocks before
implementation begins.

Unlike the Stage 3C--3G geometry and algebra checks, the full CCZ4 RHS will
not have many exact closed-form identities. A passing symbolic check for one
piece should therefore be treated as local evidence only, not as validation of
the complete evolution.

\section{Conventions}

The physical spatial dimension is
\[
  \mathrm{GR\_SPACEDIM}=4,
\]
while the Chombo/grid dimension is
\[
  \mathrm{CH\_SPACEDIM}=2.
\]
The gridded directions are \(x,z\), and the hidden transverse directions are
two equivalent \(w\) directions, so
\[
  N_{\rm hidden}
  =
  \mathrm{GR\_SPACEDIM}-\mathrm{CH\_SPACEDIM}
  =
  4-2=2.
\]
The component \(h_{ww}\) is a Cartesian-like hidden conformal metric component,
not \(g_{\theta\theta}\). The spherical/cartoon reconstruction supplies the
external radius factor,
\[
  g_{\theta\theta}
  =
  x^2\gamma_{ww}
  =
  x^2\frac{h_{ww}}{\chi}.
\]
All traces and trace-free projections are full four-dimensional spatial
operations. The conformal extrinsic-curvature split therefore uses
\[
  K_{ij}
  =
  \chi^{-1}
  \left(A_{ij}+\frac{1}{4}h_{ij}K\right).
\]

\section{Variables}

The future geometric RHS blocks are expected to involve
\[
  \chi,\quad
  h_{xx},h_{xz},h_{zz},h_{ww},\quad
  A_{xx},A_{xz},A_{zz},A_{ww},\quad
  K,\quad
  \Theta,\quad
  \hat{\Gamma}^x,\hat{\Gamma}^z,
\]
plus lapse, shift, and gauge-driver variables as needed. Stage 3H does not
choose final C++ enum names or storage order.

The connection-sector variable should use the GRChombo-facing hatted
conformal connection,
\[
  \hat{\Gamma}^i
  =
  \tilde{\Gamma}^i+2Z^i.
\]
Thus the future implementation blueprint should evolve
\(\hat{\Gamma}^A\), \(A\in\{x,z\}\), with \(Z^i\) understood as the encoded
Z4 constraint vector inside \(\hat{\Gamma}^i\). Language based directly on
\(Z^i\) remains useful for constraint-propagation analysis, but the
\(\hat{\Gamma}^i\) basis avoids an extra translation layer.

\section{RHS Blocks}

The future implementation should be decomposed into the following blocks:
\begin{itemize}
  \item \(\partial_t h_{ij}\): standard conformal-metric terms, determinant
  enforcement, hidden \(h_{ww}\) evolution, and off-diagonal \(h_{xz}\) terms.
  This block mixes inherited CCZ4 structure with modified-cartoon additions.
  \item \(\partial_t\chi\): trace and gauge terms using the full
  four-dimensional trace. This block is mostly inherited but depends on the
  four-dimensional trace bookkeeping.
  \item \(\partial_t K\): trace evolution, lapse Laplacian, Ricci scalar,
  hidden-sector geometry, and CCZ4 damping terms. Hidden Ricci pieces are
  modified-cartoon additions and require reference checks.
  \item \(\partial_t A_{ij}\): trace-free Ricci and Hessian blocks, nonlinear
  \(A\)-terms, \(A_{xz}\), and hidden \(A_{ww}\). This is a mixed block with
  high risk of silently omitted cartoon-only terms.
  \item \(\partial_t\Theta\): Hamiltonian-constraint coupling,
  constraint-propagation terms, and damping.
  \item \(\partial_t\hat{\Gamma}^A\): hatted conformal connection evolution,
  encoded \(Z^i\) contributions, momentum-constraint coupling, connection
  terms, hidden-direction divergences, and off-diagonal \(x,z\) terms.
  \item Gauge RHS blocks: lapse, shift, and any driver variables.
  \item Algebraic enforcement blocks: \(\det h_{4D}=1\), trace-free \(A\),
  hidden multiplicity, and regularity placeholders.
\end{itemize}

The gauge block owns the evolution equations for lapse, shift, and any
Gamma-driver auxiliary variables. The \(\hat{\Gamma}^A\) block owns the
occurrences of lapse, shift, and Gamma-driver quantities inside
\(\partial_t\hat{\Gamma}^A\). This ownership split should prevent
double-counting when the blocks are implemented or compared against reference
code.

\section{Validation Routes}

The future checks should be layered:
\begin{itemize}
  \item Flat or Minkowski data should give a vanishing RHS or known pure-gauge
  behavior.
  \item The sheared-flat off-diagonal Ricci gate should protect geometry-source
  handling, but it is not a full CCZ4 check.
  \item The diagonal limit \(h_{xz}=A_{xz}=\gamma_{xz}=K_{xz}=0\) should recover
  the diagonal cartoon formulas.
  \item Exact unperturbed uniform black-string data should produce a stationary
  or pure-gauge RHS once gauge conventions are fixed; the gauge-fixed uniform
  GP-string stationarity and constraint preservation outside turduckening are
  supporting anchors.
  \item Hamiltonian, momentum, \(\Theta\), and the \(Z^i\) vector encoded in
  \(\hat{\Gamma}^i\) should check constraint preservation.
  \item Symbolic dimensional extraction of the \(\Theta\), \(Z^i\), and
  \(\hat{\Gamma}^i\) terms should guard hidden multiplicity and the
  four-dimensional trace denominator.
  \item The eventual linearized RHS should reproduce the Gregory--Laflamme
  threshold/growth spectrum for the SO(3)-symmetric sector after matching
  radius convention, \(z\)-periodic boundary condition, perturbation sector,
  gauge choice, and extraction variable.
  \item Separate linearized constraint-violation injection tests around flat or
  uniform-string data should track \(\Theta\), \(Z^i\), Hamiltonian constraint,
  and momentum constraint response, including \(\kappa_1\), \(\kappa_2\),
  hidden multiplicity, and the \(1/4\) trace bookkeeping.
  \item Nonlinear source blocks require later comparison against an
  external/reference implementation and convergence testing after project
  conventions are fixed.
\end{itemize}

The Gregory--Laflamme spectrum is a semi-analytic or literature/numerical
spectral anchor from the linear GL eigenvalue problem, not a closed-form
elementary identity. Since the physical GL mode is expected to be
constraint-satisfying in the continuum, it is a strong check of the physical
linearized RHS, \(K\), \(A_{ij}\), metric, gauge, and source-sector couplings.
It is not by itself sufficient for the \(\Theta/Z^i\) damping signs, so a
separate constraint-violation damping test is required.

\section{Risk Classification}

Determinants, inverse metrics, trace-free projections, denominator factors, and
diagonal/off-diagonal limits have exact algebraic anchors. Ricci-source terms
have useful flat and sheared-flat geometry gates. Full nonlinear CCZ4 damping,
constraint propagation, gauge-driver details, enforcement timing, and late-time
GL dynamics require external/reference-code comparison and convergence evidence.

\section{Non-Goals And Next Stages}

Stage 3H does not solve small-\(x\) regularity, decide final gauge choices,
replace later external/reference comparison, or start implementation. Stage 3I
is reserved for small-\(x\) regularity and turduckening treatment. Stage 3J is
reserved for unit-test fixture design before C++ source terms are added.

\end{document}
